% Generated draft: requires mathematical and source review.
\documentclass[11pt]{article}
\usepackage[margin=25mm]{geometry}
\usepackage{fontspec}
\setmainfont{Latin Modern Roman}
\usepackage{amsmath,amssymb,mathtools,unicode-math}
\setmathfont{Latin Modern Math}
\usepackage{enumitem,needspace,xurl,hyperref}
\setlength{\parindent}{0pt}
\setlength{\parskip}{5pt}
\setlength{\emergencystretch}{3em}
\newcommand{\N}{\mathbb{N}}
\newcommand{\Z}{\mathbb{Z}}
\newcommand{\Q}{\mathbb{Q}}
\newcommand{\R}{\mathbb{R}}
\newcommand{\C}{\mathbb{C}}
\newcommand{\F}{\mathbb{F}}
\newcommand{\A}{\mathbb{A}}
\newcommand{\PP}{\mathbb{P}}
\newcommand{\E}{\mathbb{E}}
\newcommand{\eps}{\varepsilon}
\newcommand{\dd}{\,\mathrm{d}}
\begin{document}
\section*{UnsolvedMath ID: 30003161}
\textbf{Source problem code:} catalogue entry
\subsection{Definitions and mathematical statement}
The source asks for stability with respect to the flux for entropy solutions of conservation-law initial-boundary-value problems in either of two settings: several space dimensions, or one space dimension with a space-dependent flux. It does not specify enough hypotheses to determine a unique theorem.

For notation only, a general conservation-law template is
\[
\partial_t u+\operatorname{div}_x F(x,u)=0
\quad\text{in }(0,T)\times\Omega,
\qquad
u(0,\cdot)=u_0,
\]
where the initial condition is to be read with \(u\), namely \(u(0,\cdot)=u_0\), \(T>0\), \(\Omega\subseteq\R^n\) is an open spatial domain, \(u:(0,T)\times\Omega\to\R^m\), and \(F:\Omega\times\R^m\to\R^{m\times n}\) is a flux. Here
\[
(\operatorname{div}_x F(x,u))_i
=\sum_{j=1}^{n}\partial_{x_j}F_{ij}(x,u),
\qquad i=1,\ldots,m.
\]
This template must be supplemented by initial data, boundary data on \((0,T)\times\partial\Omega\), and a specified entropy and boundary admissibility formulation. Boundary data are not asserted here to hold as a pointwise Dirichlet condition. The notation permits the scalar case \(m=1\) and systems \(m>1\); the supplied statement does not determine which is intended. Likewise, writing \(F(x,u)\) allows spatial dependence without asserting that it is required in the multidimensional case.

\textbf{Requested stability problem.} Establish a mathematically specified form of continuous dependence of entropy solutions on the flux for the corresponding initial-boundary-value problems, in the multidimensional setting \(n\geq2\), or even in the setting \(n=1\) with \(F\) depending on \(x\). A schematic continuity target, once an admissible problem class and convergence notions have been fixed, is
\[
F_k\longrightarrow F
\quad\Longrightarrow\quad
u_k\longrightarrow u,
\]
where \(u_k\), rather than \(\nu_k\), denotes the entropy solution associated with \(F_k\), and \(u\) denotes the solution associated with \(F\). Thus the intended schematic implication is \(F_k\to F\Longrightarrow u_k\to u\). Comparing problems with common initial and boundary data is one possible formulation of dependence on the flux; the source does not specify whether data are fixed or allowed to vary.

\textbf{Unresolved qualifications.} The admissible fluxes, scalar or system scope, spatial domains, data classes, entropy conditions, boundary interpretation, solution spaces, flux and solution topologies, and any quantitative stability estimate are unspecified. The schematic target therefore records the requested direction of dependence and is not a fully quantified conjecture. No universal claim over all domains, fluxes, or entropy formulations is intended.
\subsection{Short English statement}
The problem asks whether changing the flux slightly changes entropy solutions only slightly in conservation-law initial-boundary-value problems, either in several space dimensions or in one dimension with spatially varying flux. The source leaves the admissible equations, boundary formulation, and meaning of stability unspecified, so a precise theorem requires further hypotheses.
\subsection{Sources}
\begin{itemize}
\item Local registry: \texttt{data/unsolvedmath/problems.json}, supplied record with permanent UnsolvedMath ID \(30003161\), source problem code \texttt{catalogue entry}, title \emph{Flux Stability for Multidimensional Conservation Laws}; collection: Oberwolfach Reports Open Problems. No individual proposer is supplied.
\item Supplied report citation: \emph{Hyperbolic Techniques in Modelling, Analysis and Numerics} (2017), Oberwolfach Reports, DOI \texttt{10.4171/owr/2016/30}, pp.~1683--1751. \url{https://doi.org/10.4171/owr/2016/30}. The registry separately records proposed year 2016; the differing year labels are retained without reconciliation.
\item Supplied literature source: Oberwolfach Report 30. \url{https://ems.press/content/serial-article-files/46634?nt=1}.
\item Supplied literature source: Rossi, stability for scalar conservation laws with boundaries. \url{https://arxiv.org/abs/1809.06066}. This is the supplied descriptive attribution; full bibliographic details are not provided.
\item Literature assessment dated exactly \texttt{2026-08-21}: the supplied curation labels the problem open and states that stability results are known for important one-dimensional scalar settings, while the requested extension to multidimensional or space-dependent-flux initial-boundary problems remains open. This is supplied dated curation derived from the OWR record, not an independently verified assessment of the literature or present mathematical status.
\item The registry labels the statement \texttt{corrected\_verified}, with the qualification that notation, typography, and sentence structure were repaired without changing mathematical content. This supplied editorial label does not establish mathematical correctness or independent source verification.
\end{itemize}
\end{document}
